%================================
% Vocabulary & Language Focus — XB1 (D–F)
%================================
\documentclass[../main.tex]{subfiles}

\begin{document}

%================================
\section{D}
%================================

\vw{damp}{/dæmp/}{adj.}

\sense{adj.}{潮湿的，湿气重的}
\col{damp clothes / walls / weather}
\ex{Don't sleep between damp sheets --- you'll catch a cold.}
\ex{The walls of the basement were cold and damp.}
\col{damp sth down（抑制；封火）}
\ex{She damped down her anger and forced a smile.}

\begin{lex}{damp / wet / moist / humid}
湿度递进看对象：\textbf{wet} 最普通，“湿的、未干的”；\textbf{moist} 微湿且常含好感（\textit{moist cake} 松软湿润的蛋糕）；\textbf{damp} 令人不适的潮气（\textit{a damp room}）；\textbf{humid} 专指空气湿度大（\textit{a humid climate}）。另外 damp 可作名词“潮气”：\textit{The damp in the house aggravated his cough.}
\end{lex}

\begin{origin}
14 世纪借自古诺斯语 \textit{dampi}“\textbf{尘埃、水汽}”及 \textit{demma}“使阻遏”（与德语 \textbf{Dampf}“蒸汽”、\textit{dämpfen}“抑制、使消音”同源，原始日耳曼语 *\textit{dampaz}）。德语物理学的 \textbf{Dämpfung}（阻尼）即此词——工程英语 \textit{dampen}“使减振”保留原义，而 \textit{dampen one's enthusiasm}“泼冷水”是它最生动的引申：把蒸汽压回去。
\end{origin}

\vw{dare}{/deər/}{vi. \& vt. \& mod.\ v.}

\sense{vi. \& vt.}{敢于，胆敢}
\col{dare (to) do sth（实义动词用法）}
\ex{He didn't dare (to) tell his parents the truth.}
\ex{She would never dare to speak to the boss like that.}
\col{not dare (to) do sth / daren't do sth（情态动词用法，无 to、无第三人称 -s）}
\ex{He daren't say a word in case he makes it worse.}
\ex{How dare you talk to me like that!}

\sense{vt.}{激（某人做某事）}
\col{dare sb to do sth（挑激对方去做）}
\ex{They dared him to jump off the bridge into the river.}
\col{I dare say（我敢说，英式插入语）}
\ex{I dare say you're right, but we still need proof.}

\begin{gram}[dare：情态动词与实义动词的双重身份]
dare 与 need 是英语中仅有的两个\textbf{“半情态动词”}，两套语法并存：
\begin{center}
\begin{tabular}{lll}
\toprule
 & 情态用法 & 实义用法 \\
\midrule
否定 & \textit{He \textbf{daren't} go.} & \textit{He \textbf{doesn't dare to} go.} \\
疑问 & \textit{\textbf{Dare} he go?} & \textit{\textbf{Does} he \textbf{dare (to)} go?} \\
\bottomrule
\end{tabular}
\end{center}
要点：（1）否定句中 \textit{don't dare to} 的 to 可省，遂产生混血形式 \textit{He doesn't dare go}——完全合法；（2）\textit{I dare say} 已固化为习语“想必、大概”；（3）\textit{How dare you + 动词原形}表愤怒斥责；（4）dare 作“激将”义时\textbf{只能是实义动词}（\textit{dare sb \textbf{to} do}）。语法填空中判断 dare 身份的线索：前面有无 do/does/did。
\end{gram}

\begin{origin}
古英语 \textit{durran}“敢于”，是动词 \textit{durran/durr}（现 \textbf{durst}，其过去式）的现在时，源自原始日耳曼语 *\textit{durzanai}，字面义“\textbf{能够下决心}”，与哥特语 \textit{gadaursan} 同源；最终可追至原始印欧语 *\textit{dhers-}“\textbf{大胆、勇猛}”。希腊语同根词 \textit{thrasys}“大胆的”（\textbf{thrason}“自大吹牛者”，英语 \textit{thrasonical}“自负的”即源于此，出自泰伦提乌斯喜剧人物 Thraso）。一个“勇敢”的词根，在希腊长成“吹牛”，在英格兰长成“敢”——词义的分化如同物种的演化。
\end{origin}

\vw{debt}{/det/}{n.}

\sense{n.}{人情债，情义，恩情}
\col{be in sb's debt / in debt to sb（欠某人情）}
\ex{I'm deeply in your debt for what you did that night.}
\col{repay a debt（还债；报恩）}
\ex{He worked for years to repay his debt to the man who saved him.}

\sense{n.}{借款，欠款，债务}
\col{be (deeply) in debt / get (out of) debt}
\ex{The company got deep into debt after two losing seasons.}
\col{pay off / run up a debt（还清／背上债务）}
\ex{She took a second job to pay off her student debts.}
\col{national debt（国债）}
\ex{The national debt has doubled over the past decade.}

\begin{gram}[哑音 b 与“词源矫正”拼写]
debt 的 /t/ 后藏着不发音的 b——这不是英语本有的，而是 15 世纪学者依据拉丁语 \textit{debitum} 人为加回去的（法语借入时本作 \textit{dette}）。同类“伪造”拼写：\textbf{doubt}（$<$拉丁 \textit{dubitare}）、\textbf{island}（b 是误仿 \textit{isle} 加的——isle 的 s 才是真词源 $<$拉丁 \textit{insula}，而 island 是纯日耳曼词 \textit{īegland}）、\textbf{receipt}、\textbf{subtle}。听音写词时千万别被这些“学术装饰音”骗了。
\end{gram}

\begin{origin}
古法语 \textit{dette} $\leftarrow$ 拉丁语 \textit{debitum}“\textbf{所欠之物}，\textit{debēre}“欠”的过去分词中性：\textit{de-}（除去、离开）+ \textit{habēre}“\textbf{有}”——“欠”就是“从你该有的当中\textbf{拿走了}一部分”。这个词根组合在英语中极为高产：\textbf{debit}（借方）、\textbf{due}（应付的，\textit{debēre} 的另一后代）、\textbf{duty}（职责“本应尽的本分”）、\textbf{debt} 的兄弟 \textbf{debenture}（债券）。双关老话 \textit{``Debt is the other man's dollar.''} 也植根于“de-(不) + habere(拥有)”的窘境。
\end{origin}

\vw{decade}{/ˈdekeɪd/}{n.}

\sense{n.}{十年，十年期}
\col{over / in / during the past decade}
\ex{The city has changed beyond recognition over the past decade.}
\col{decades of\dots（数十年）}
\ex{It took decades of research to confirm the theory.}
\col{for decades on end（连续几十年）}
\ex{The lighthouse keeper had worked there for decades on end.}

\begin{origin}
法语 \textit{décade} $\leftarrow$ 希腊语 \textit{dekas (dekados)}“\textbf{十的一组}”，出自 \textit{deka}“十”（原始印欧语 *\textit{dekm̥-}“十”——这个基数词在各语族惊人地同源：十、\textit{decem}（拉）、\textit{dasa}（梵）、\textit{zehn}（德）。同根词 \textbf{December}（罗马旧历的十月）、\textbf{decathlon}（十项全能）、\textbf{decimal}（十进制的）、\textbf{dime}（美式一角，$<$古法语 \textit{disme}“十分之一”）。顺便一提：英语“十”\textbf{ten} 亦出自同一原始印欧词——decade 与 ten 是隔着六千年时空的孪生兄弟。
\end{origin}

\vw{decline}{/dɪˈklaɪn/}{vi.\&vt.\ \&\ n.}

\sense{vi. \& vt.}{减少，下降，衰退，衰落}
\col{decline by\dots（下降若干幅度）}
\ex{Sales declined by 15 percent in the fourth quarter.}
\col{a declining population / market（渐减的）}
\ex{The school struggles with a declining number of students.}

\sense{vi. \& vt.}{谢绝（正式礼貌的拒绝）}
\col{decline an invitation / an offer}
\ex{She declined the invitation, citing a prior engagement.}
\col{decline to do sth}
\ex{The spokesman declined to comment on the rumors.}

\sense{n.}{衰退，衰落，减少，下降}
\col{a decline in sth / fall into decline / in decline}
\ex{There has been a sharp decline in reading among teenagers.}
\ex{The old port fell into decline when the shipping moved elsewhere.}

\begin{lex}{decline / refuse / reject}
\textbf{decline} 最委婉正式，“婉言谢绝”，主语通常是人；\textbf{refuse} 直截了当“拒绝”，可接名词或不定式（\textit{refuse to obey}）；\textbf{reject} 语气最强，多指\textbf{否决、驳回、拒收}（\textit{reject a proposal}），带“不认可”的判断色彩。写作中婉拒对方来信邀约，宜用 decline。
\end{lex}

\begin{origin}
拉丁语 \textit{declinare}“\textbf{使偏斜、转开}”：\textit{de-}（向下、离开）+ \textit{clinare}“\textbf{倾斜}”（原始印欧语 *\textit{klei-}“ lean 倾侧”，同根词 \textbf{lean}（倾斜$<$古英语 \textit{hlīnan}，实为拉丁词的日耳曼表亲）、\textbf{client}（*\textit{klient-}“倚靠别人的人”$\to$依附者$\to$客户）、\textbf{climate}（希腊 \textit{klinein}“倾斜”——古人认为地表的“斜度”决定气候）、\textbf{clinic}（\textit{klinikē}“卧床的医学”）。“谢绝”即“把头\textbf{偏向一边}”婉拒；语法术语 \textbf{declension}（变格）“名词的倾斜变化”同源。
\end{origin}

\vw{defend}{/dɪˈfend/}{vt. \& vi.}

\sense{vt. \& vi.}{保护，防御}
\col{defend sb/sth from/against sb/sth}
\ex{The castle walls were built to defend the city against invaders.}
\ex{A mother will defend her young with her life.}

\sense{vt.}{辩白，辩解}
\col{defend oneself（自卫；为自己辩护）}
\ex{The minister defended herself against the charges in Parliament.}
\col{defend one's thesis（答辩）}
\ex{She successfully defended her doctoral thesis last week.}

\begin{gram}{defend 的反义家族：fend 家族的双向演变}
拉丁语 \textit{fendere}“\textbf{击打}”加前缀后成对出现：\textbf{de-fend}（打下去$\to$防御）、\textbf{of-fend}（打过去$\to$冒犯；offence 犯罪；进攻\/防守在球场上是一枚硬币的两面，故 \textit{offense} 美式英语兼指“进攻”）、\textbf{de-fence}（defend 的名词，又名词转动词“围以栅栏”，\textit{fence} 篱笆、击剑、销赃者 fence 均由此出）。认准词根 \textit{fend/fens}，一串词全部拿下：\textbf{defensive, defendant（被告“辩护者”）, offensive, fencing}。
\end{gram}

\begin{origin}
拉丁语 \textit{defendere}“\textbf{击退、挡开}”：\textit{de-}（离开）+ \textit{fendere}“敲、刺”（词源不甚明朗，一说与原始印欧语 *\textit{gwhend-}“击、抓”有关，与英语 \textbf{beat} 无关但功能相当）。中世纪拉丁语 \textit{defensa}“守护的”给了法语 \textit{défense}。法律用语 \textbf{defendant}“被告”字面即“为自己辩护的人”；足球的 \textbf{defender} 后卫与 \textbf{defense line} 防线皆是“挡开攻击者”的直译。
\end{origin}

\vw{delicate}{/ˈdelɪkət/}{adj.}

\sense{adj.}{精致的，精细的，精密的}
\col{delicate craftsmanship / a delicate instrument}
\ex{The watch contains some of the most delicate mechanisms ever made.}
\ex{The vase was painted with delicate brushstrokes.}

\sense{adj.}{易损的，易碎的，脆弱的}
\col{delicate skin / fabrics / ecosystems}
\ex{The coral reef is a delicate ecosystem easily damaged by pollution.}
\col{in delicate health（体弱多病的，委婉语）}
\ex{Grandma has been in delicate health since winter.}
\col{a delicate situation / matter / issue（需小心处理的）}
\ex{The negotiators trod carefully through a series of delicate issues.}

\begin{origin}
拉丁语 \textit{delicatus}“\textbf{令人愉悦的、娇美的}”，源自 \textit{delicia}“快乐、爱物”（\textit{deliciae} 复数“心上人、宠儿”），动词 \textit{delicere}“\textbf{引诱}”：\textit{de-} + \textit{lacere}“诱惑、设饵”。\textit{lacere} 正是法语 \textbf{lacer}“勾住”与英语 \textbf{delight}（$<$古法语 \textit{delitier} $\leftarrow$ \textit{delictare}“使愉悦”，delectable 亦同）的祖先——\textbf{delicate} 与 \textbf{delight} 是同一个“诱人”词根的两支后代。语义从“魅惑人心的精美”引申至“精美易损的”，恰如精致之物往往脆弱。
\end{origin}

\vw{depress}{/dɪˈpres/}{vt.}

\sense{vt.}{使抑郁，使沮丧}
\ex{The gloomy weather depresses everyone in the office.}
\col{be / feel depressed（沮丧的；抑郁症的）}
\ex{He had been depressed for months after losing his job.}

\sense{vt.}{使萧条，使不景气}
\col{depress prices / the market（压低）}
\ex{Cheap imports depressed the price of local rice.}
\ex{The tourism industry was depressed by the travel restrictions.}

\begin{gram}{prem-“压”词族：一组前缀的全家福}
拉丁语 \textit{premere}“压”是英语最整饬的词族之一：\textbf{press}（压）、\textbf{de-press}（向下压$\to$使沮丧、使萧条）、\textbf{com-press}（压在一起$\to$压缩）、\textbf{ex-press}（压出来$\to$挤出言辞$\to$表达；快车“挤出时间直达”）、\textbf{im-press}（压入$\to$印；使铭记$\to$impressive）、\textbf{re-press}（压回去$\to$镇压；压抑情感）、\textbf{sup-press}（sub- 压到底$\to$查禁；忍住）。名词形式一律 -pression：depression（抑郁；大萧条）、expression、impression。一词根贯穿心理学（depression）、经济学（the Great Depression）与印刷史（press 印刷机$\to$新闻界）。
\end{gram}

\vw{deserve}{/dɪˈzɜːv/}{vt.}

\sense{vt.}{值得，应得}
\col{deserve sth / deserve to do sth}
\ex{She deserves a long holiday after all that hard work.}
\ex{They deserved to win --- they played brilliantly.}
\col{deserve doing / deserve to be done（主动动名词表被动！)}
\ex{These young criminals deserve punishing.（= deserve to be punished）}
\col{get what one deserves / serving sb right}
\ex{He got exactly what he deserved; it served him right.}

\begin{gram}{deserve / need / want + doing：动名词主动式表被动义}
\textit{The proposal deserves \textbf{considering}.}（= deserves to be considered）——在 deserve, need, want, require 后，动名词的\textbf{主动形式}表达\textbf{被动意义}：\textit{The car needs cleaning.}（车需要被洗）。\ding{55} \textit{needs being cleaned}。这与 be to blame 主动表被动（见 blame 词条）同属一种经济原则：当主语是“该受动作影响者”时，英语倾向省去被动标记。高考语法填空的高频陷阱。
\end{gram}

\begin{origin}
拉丁语 \textit{deservire}“\textbf{殷勤地侍奉、效忠}”：\textit{de-}（彻底地，加强语气）+ \textit{servire}“服务”（\textit{servus}“奴仆”，同根 \textbf{serve, service, servant, conservation}“保持原样的服务”、\textbf{desert}“应得的赏罚之地”$<$\textit{desertum}，注意区别于 \textbf{desert} 沙漠 $<$\textit{desertare}“遗弃”）。古罗马的“应得”以\textbf{服务}为前提：先彻底效力，才有资格论功行赏。习语 \textit{get one's just deserts}“得到应有的下场”保留了这个词的古拼写与古义。
\end{origin}

\vw{desperate}{/ˈdespərət/}{adj.}

\sense{adj.}{绝望的}
\ex{The trapped miners were getting desperate after five days underground.}

\sense{adj.}{不惜冒险的（孤注一掷的）}
\col{a desperate attempt / measure}
\ex{In a desperate attempt to escape, the fox gnawed off its own leg.}

\sense{adj.}{极其需要的}
\col{be desperate for sth / to do sth}
\ex{The refugees were desperate for clean water and medicine.}
\ex{She was desperate to prove herself to the coach.}

\sense{adj.}{极严重的（局面）}
\col{a desperate shortage / situation}
\ex{The hospital faces a desperate shortage of nurses.}

\begin{origin}
拉丁语 \textit{desperatus}“\textbf{彻底无望}”，\textit{desperare}“绝望”的过去分词：\textit{de-}（去除）+ \textit{sperare}“\textbf{希望}”（名词 \textit{spes}“希望”，源自原始印欧语 *\textit{spē-}“繁荣、成功”）。同根词 \textbf{despair}（名词“绝望”，desperate 的名词近亲）、\textbf{prosper}（pro- 向着希望$\to$兴旺；\textit{prosperity}）、以及出人意料的 \textbf{speed}（古义“成功、兴隆”，\textit{God speed you!}“愿神使你发达”$\to$\textit{speeding} 祝顺利；\textit{more haste, less speed} 中它仍指“顺利”）。“绝望”即“希望的彻底除名”，与汉语“心如死灰”取象相通。
\end{origin}

\vw{detect}{/dɪˈtekt/}{vt.}

\sense{vt.}{发现，查明，侦察出}
\col{detect sth in sb's voice / behavior（察觉）}
\ex{I detected a note of sarcasm in her otherwise friendly voice.}
\col{be detected by radar / sensors}
\ex{Small drones can be detected by radar at a distance of ten kilometers.}
\col{detection / detector / detective（名词族）}
\ex{Smoke detectors should be checked every month.}
\ex{The detective gathered evidence for weeks before making an arrest.}

\begin{origin}
拉丁语 \textit{detegere}“\textbf{揭开盖子}”：\textit{de-}（除去）+ \textit{tegere}“\textbf{覆盖}”（原始印欧语 *\textit{(s)teg-}“盖”，同根 \textbf{thatch} 茅草屋顶、\textbf{deck} 甲板“铺板”、\textbf{tile} 瓦、\textbf{integument} 皮肤“覆盖层”）。“侦查”就是把真相从遮盖物下\textbf{揭出来}。反义词 \textbf{de-tegere} 的对立面 \textit{tegere} 加前缀 \textit{pro-}（在前）即是 \textbf{protect}“挡在前面$\to$保护”——detect 与 protect 是同一块“盖子”的两面：一个揭开，一个盖上。
\end{origin}

\vp{dig up}

\sense{v. 短语}{发现，搜集，查明（新闻、丑闻）}
\ex{Reporters dug up some old photos the senator hoped nobody would see.}
\ex{Investigators spent months digging up evidence of the fraud.}
\col{dig sth up（本义：挖出，掘出）}
\ex{Workers digging up the road found an ancient bronze vessel.}

\vw{digest}{/daɪˈdʒest/ vt.\ \&\ vi.\ \ /ˈdaɪdʒest/ n.}{vt.\&vi.\ \&\ n.}

\sense{vt. \& vi.}{消化}
\ex{Fatty food takes much longer to digest than vegetables.}
\col{be hard/easy to digest（难／易消化；喻“难理解”）}
\ex{The news was hard to digest all at once.}

\sense{vt.}{领会，理解}
\ex{It took me a while to digest the implications of her words.}

\sense{n. /ˈdaɪdʒest/}{摘要，文摘}
\ex{\textit{Reader's Digest} was one of the best-selling magazines of the century.}

\begin{origin}
拉丁语 \textit{digerere}（过去分词 \textit{digestus}）“\textbf{分开携带、疏散、消化}”：\textit{di-}（dis- 分开）+ \textit{gerere}“\textbf{携带、进行、做}”。\textit{gerere} 是个多产的“办事”动词：\textbf{gesture}（做出来的姿势）、\textbf{ex-gerere}（带出来$\to$\textbf{exaggerate} 夸张“带过了头”）、\textbf{be-gerere}（\textbf{belligerent} 好战的“扛着战争的人”）、\textbf{gerund}（“该做的事”$\to$动名词，语法学家的玩笑）。“消化”的本义是“把食物\textbf{分装}运走”；“文摘”是罗马法学家对浩瀚典籍的“分装提要”——查士丁尼皇帝的《学说汇纂》即名 \textit{Digesta}。注意名动同音 /ˈdaɪdʒest/，是“重音转移构词法”（见 advocate 词条）的又一例。
\end{origin}

\vw{dimension}{/daɪˈmenʃn/}{n.}

\sense{n.}{方面，侧面；规模，程度}
\col{a new / different dimension（新的层面）}
\ex{The Internet has added a whole new dimension to shopping.}
\col{the dimensions of the problem（问题的规模）}
\ex{The full dimensions of the scandal only became clear last month.}

\sense{n.}{维（空间）}
\col{three-dimensional / in three dimensions}
\ex{The movie was praised for its vivid three-dimensional effects.}
\col{measure the dimensions of a room（长宽高尺寸）}
\ex{First measure the dimensions of the room before buying furniture.}

\begin{origin}
拉丁语 \textit{dimensio}（\textit{dimetiri} 的名词）“\textbf{量出}”：\textit{dis-}（分开）+ \textit{metiri}“\textbf{测量}”（原始印欧语 *\textit{mē-}“测量”，同根的浩大词族：\textbf{measure}、\textbf{meter}（米；计量器）、\textbf{metric}、\textbf{immense}（im- 无法测量$\to$无边无际）、\textbf{commensurate}（可同量$\to$相称的）、\textbf{dimension} 的兄弟 \textbf{premeditated}（预先掂量$\to$预谋的）。“维”的数学义出自“量出的方向”。“俚语亲戚”还有 \textbf{meal}（古义“时间、次数”，一日\textbf{几次}饭）——测量的词根竟藏在一日三餐里。
\end{origin}

\vw{display}{/dɪˈspleɪ/}{vt.\ \&\ n.}

\sense{vt.}{陈列，展出；显示，表现}
\col{display sth in a shop / museum}
\ex{The museum displays porcelain from the Song dynasty.}
\col{display one's talent / courage / emotions}
\ex{She displayed remarkable courage during the earthquake.}

\sense{n.}{陈列，展览；表现；展示}
\col{on display（在展出中，高频）}
\ex{Some of her early paintings are on display at the city gallery.}
\col{a firework / flying display（表演）}
\ex{Thousands watched the spectacular flying display.}
\col{a display of affection / wealth（……的展示）}
\ex{He never made any display of his wealth.}

\begin{origin}
经由古法语 \textit{despleier} $\leftarrow$ 晚期拉丁语 \textit{displicare}“\textbf{展开、摊开}”：\textit{dis-}（分开）+ \textit{plicare}“\textbf{折叠}”（原始印欧语 *\textit{plek-}，即 complex 词条的“编织”词根的近亲）。\textit{plicare} 的后代们方向各异：折起来是 \textbf{ply}（折叠；层），展开摊平是 \textbf{display} 与 \textbf{deploy}（军队“展开”队形$\to$部署），缠绕进去是 \textbf{employ/imply/implicate}（\textit{im-plicare}“折入其中”$\to$卷入），折得复杂是 \textbf{complicated}。认准 \textit{plic-/ply}“折”，一整页词汇表应声而解。
\end{origin}

\vw{distant}{/ˈdɪstənt/}{adj.}

\sense{adj.}{遥远的，远处的，久远的}
\col{a distant village / country / memory}
\ex{The thunder of the falls could be heard from a distant valley.}
\col{in the distant past / future（久远的）}
\ex{In the not-too-distant future, most cars may drive themselves.}

\sense{adj.}{冷淡的}
\ex{After the quarrel, she remained cold and distant for weeks.}

\sense{adj.}{远亲的}
\col{a distant relative / cousin}
\ex{We discovered we were distant cousins through a DNA test.}

\begin{origin}
拉丁语 \textit{distare}“\textbf{站立相隔}”：\textit{di(s)-}（分开）+ \textit{stare}“站”（\textit{stā-} 家族，见 consist 词条——“遥远的”就是“站在很远处”）。\textbf{distance}（距离）是它的名词，\textbf{distant} 是形容词；同一家族的“分开站”还有 \textbf{distinguish}（见下条）与 \textbf{estate}（\textit{ex-stare}“立在外面”$\to$地产、状态）。社交上的“冷淡疏远”是空间隐喻的直接延伸，心理学称之为“物理距离的社交转喻”，cold shoulder（冷肩膀）与 distant 在各语言中都不谋而合。
\end{origin}

\vw{distinguish}{/dɪˈstɪŋɡwɪʃ/}{vt. \& vi.}

\sense{vt. \& vi.}{区分，辨别；认出}
\col{distinguish A from B / distinguish between A and B}
\ex{Children often cannot distinguish right from wrong without guidance.}
\ex{It is hard to distinguish between the twin sisters at a glance.}

\sense{vt.}{成为……的特征，使有别于}
\col{be distinguished by / what distinguishes A from B is\dots}
\ex{What distinguishes our school is its emphasis on independent thinking.}

\sense{vt.}{使出众}
\col{distinguish oneself（使自己出类拔萃）}
\ex{He distinguished himself as a fearless battlefield surgeon.}
\col{a distinguished scientist / career（卓著的）}
\ex{A distinguished chemist delivered the opening lecture.}

\begin{gram}{“区分”句式的三种写法与写作应用}
同一逻辑，三种句式：（1）\textit{distinguish A \textbf{from} B}；（2）\textit{distinguish \textbf{between} A and B}；（3）被动式 \textit{A is distinguished from B \textbf{by} C}（以 C 区别于）。在比较类作文中，（3）是漂亮的主题句模板：\textit{Genuine friendship is distinguished from mere acquaintance by mutual trust.} 注意 distinguished 与 distinguishable（可辨别的）词义不同，勿混。
\end{gram}

\begin{origin}
拉丁语 \textit{distinguere}“\textbf{用刺记分开}”：\textit{di(s)-}（分开）+ \textit{stinguere}“刺、戳”（词源上与英语 \textbf{sting} 同属原始印欧语 *\textit{steg-}“尖刺”家族）。古罗马牧人以刺标（在牲口身上刺记号）区分自家的羊——“辨别”的本义是“打上刺印分开”。同根的 \textbf{extinguish}（ex- 熄灭“把火刺灭”？实为“扑灭、除掉”，词根相同）与 \textbf{distinct}（可分辨开的）、\textbf{instinct}（in- 内在的刺$\to$本能“扎在体内的冲动”）。
\end{origin}

\vw{district}{/ˈdɪstrɪkt/}{n.}

\sense{n.}{区，行政区；地区，区域}
\col{a business / residential / historic district}
\ex{The old town has been preserved as a historic district.}
\col{a school / police / electoral district（管辖或选区的“区”）}
\ex{Each school district sets its own holiday calendar.}

\begin{lex}{district / area / region / zone}
\textbf{district} 带\textbf{行政或功能}色彩（行政区、管区、特色街区）；\textbf{area} 最泛的“地区、面积”；\textbf{region} 大范围的地理区域（\textit{the Arctic region}）；\textbf{zone} 人为划定的专门地带（\textit{a war zone} 战区、\textit{a pedestrian zone} 步行区、时区 \textit{time zone}）。写作选词须先问：这一“区”是行政的、地理的，还是划定的？
\end{lex}

\begin{origin}
13 世纪借自法语，源自拉丁语 \textit{distringere}“\textbf{拉开、拽住、羁押}”：\textit{di(s)-} + \textit{stringere}“\textbf{捆紧}”（\textbf{strict}、\textbf{strain}、\textbf{stress}、\textbf{stringent} 严格的、\textbf{constrain} 约束皆此词根）。中世纪的 \textit{districtus} 指\textbf{封建领主司法权所及的范围}——“把人拽得到”的地界；故它与 \textbf{distinguish} 表面相似、词根却不同（那边的 -stingu- 是“刺”，这边的 -string- 是“捆”）。一个词，藏着半部中世纪法律地理史。
\end{origin}

\vp{do the trick}

\sense{v. 短语}{奏效，起作用，达到目的}
\ex{I had a headache, but a cup of strong tea did the trick.}
\ex{Just restart the router --- that usually does the trick.}

\begin{gram}{“trick”短语一家}
\textbf{do the trick} 奏效（trick 此处取 17 世纪“诀窍、把戏”义）；\textbf{the trick of the trade} 行家的诀窍；\textbf{trick sb into doing sth} 哄骗某人做某事；\textbf{trick or treat} 不给糖就捣乱（万圣节，treat“款待”与 trick 押头韵）；\textbf{not miss a trick} 明察秋毫；\textbf{a trick of the light} 光线造成的错觉。惯用法中 trick 多带中性“手段”义，并非全然贬义。
\end{gram}

\vw{dominate}{/ˈdɒmɪneɪt/}{vt. \& vi.}

\sense{vt. \& vi.}{在……中具有最重要（或明显）的特色；占有优势}
\col{dominate the market / headlines / the match}
\ex{Two companies dominate the market for electric cars.}
\ex{The issue of housing dominated the headlines all summer.}

\sense{vt.}{支配，控制}
\ex{A tough old man dominated the family for decades.}
\col{be dominated by（被……主导）}
\ex{The discussion was dominated by questions about safety.}

\sense{vi. \& vt.}{俯视（高地、建筑）}
\ex{A huge statue of the general dominates the town square.}

\begin{origin}
拉丁语 \textit{dominari}“\textbf{作主人、统治}”，源自 \textit{dominus}“\textbf{主人、领主}”——即 \textit{domus}“\textbf{房屋}”的“户主”。\textit{domus}（原始印欧语 *\textit{dem-}“建造房屋”）是人类最古老的词汇之一，其后代遍布欧洲：\textbf{dome}（圆顶屋）、\textbf{domestic}（家里的$\to$驯化的；国内的）、\textbf{domain}（领地；领域）、\textbf{dame / damsel}（女主／少女）、\textbf{predominant}（pre- 在前的$\to$占主导的）、甚至 \textbf{danger}（$<$\textit{dominium}“领主的权力”$\to$“受制于权势的风险”）。“统治”的原型不过是“\textbf{当家}”——从 domus 到 dominate，浓缩了一部家庭权力放大为国家权力的历史。
\end{origin}

\vw{dramatic}{/drəˈmætɪk/}{adj.}

\sense{adj.}{激动人心的，引人注目的；突然的}
\col{a dramatic increase / change / improvement}
\ex{The new policy brought a dramatic rise in exports.}
\ex{The rescue was one of the most dramatic moments of the night.}

\sense{adj.}{戏剧的}
\col{dramatic irony / a dramatic monologue}
\ex{She teaches dramatic literature at the conservatory.}

\sense{adj.}{夸张的（贬义，口语）}
\ex{He's always been a little dramatic --- every small problem becomes a tragedy.}
\col{dramatically（急剧地）}
\ex{Car ownership rose dramatically between 2000 and 2020.}

\begin{origin}
经由晚期拉丁语 \textit{dramaticus} $\leftarrow$ 希腊语 \textit{dramatikos}，出自 \textit{drama}“\textbf{做、行动、上演之事}”（动词 \textit{dran}“\textbf{做、干”}，与原始印欧语 *\textit{der-}“做、工作”相关）。与 \textit{dran} 同根的还有 \textbf{dromedary}（单峰驼“奔跑的骆驼”，\textit{dromos}“跑”）、\textbf{hippodrome}（\textit{hippos} 马 + \textit{dromos} 跑道$\to$竞技场）、\textbf{syndrome}（\textit{syn-} 一起 + \textit{dromos}“一同跑来”$\to$综合征）。故 drama 的本义是“\textbf{正在做}”——戏剧是行动的艺术，亚里士多德《诗学》开篇即如此定义。
\end{origin}

\vw{dynasty}{/ˈdɪnəsti/}{n.}

\sense{n.}{朝，代；王朝}
\col{the Ming Dynasty / a dynasty of rulers}
\ex{The porcelain was made during the Qing Dynasty.}
\col{a political / business dynasty（引申：世家的权力延续）}
\ex{The family built a business dynasty over three generations.}

\begin{origin}
经由法语 \textit{dynastie} $\leftarrow$ 晚期拉丁语 \textit{dynastia} $\leftarrow$ 希腊语 \textit{dynasteia}“\textbf{权力、霸权}”，出自 \textit{dynasthai}“\textbf{能够、有能力”}（原始印欧语 *\textit{dheh\u{1}-}“哺乳、建立”的音变形式）。同根的科学词火力全开：\textbf{dynamic}（有力的）、\textbf{dynamics}（动力学）、\textbf{dynamite}（炸药“威力之物”，诺贝尔时代的商标名）、\textbf{dynamo}（发电机）、\textbf{aerodynamics}（空气动力学）。“王朝”的定义由此清晰：不是血统，而是\textbf{能力（权力）的世代延续}。
\end{origin}

%================================
\section{E}
%================================

\vw{ease}{/iːz/}{n.\ \&\ vi.\&vt.}

\sense{n.}{容易，轻易}
\col{with ease（毫不费力地）}
\ex{The champion won the final with ease.}
\col{the ease of doing sth}
\ex{Users praise the app for its ease of operation.}

\sense{n.}{舒适，安逸}
\col{at ease / (ill) at ease（自在／局促不安）}
\ex{Put the new students at ease before the interview.}
\ex{He felt distinctly ill at ease at the formal dinner.}

\sense{vi. \& vt.}{（使）宽慰，减轻}
\col{ease the pain / pressure / burden}
\ex{A hot bath eased the pain in his shoulders.}
\col{ease one's mind（使宽心）}
\ex{A phone call home eased her mind considerably.}
\col{ease off / ease up（缓和，放松）}
\ex{The storm eased off toward midnight.}
\col{ease sb into sth（使逐渐适应）}
\ex{The coach eased the young player into the first team.}

\begin{origin}
中古英语 \textit{aise} $\leftarrow$ 古法语 \textit{aise}“\textbf{舒适、便利、机会}”，多数词源学家认为它来自通俗拉丁语 *\textit{adjacens}（\textit{adjacēre}“\textbf{躺在近旁}”的现在分词“躺卧于侧”）——即 \textbf{adjacent}（邻近的）的同源词：“舒适”就是“\textbf{躺卧在合适的近处}”。另有说法与 \textit{asiatum}“方便的处所”相关。无论哪一支，法语 \textit{aise} 都先有“身体安顿”之义，再精神化为“安心”；16 世纪英语中加上了想象出来的 h（同 \textit{host, island}），拼作 \textbf{ease}。其否定前缀形式 \textbf{dis-ease}（不安$\to$疾病）比 disease 的现代用法更见本义：病即“\textbf{不自在}”。
\end{origin}

\vw{edge}{/edʒ/}{n.}

\sense{n.}{边，边缘}
\col{on / at the edge of sth（本义与喻义）}
\ex{They pitched their tent on the edge of the forest.}
\ex{The company is on the edge of bankruptcy.}
\col{have an edge over sb（略占上风，口语高频）}
\ex{Bilingual applicants have an edge over the others.}

\sense{n.}{刀口}
\col{the edge of a blade / knife}
\ex{The knife has lost its edge and needs sharpening.}
\col{on edge（紧张不安）；take the edge off sth（减弱，缓和）}
\ex{Too much coffee makes me nervous and on edge.}
\ex{A little food took the edge off her hunger.}

\begin{origin}
古英语 \textit{ecg}“\textbf{刀刃、锋”}，源自原始日耳曼语 *\textit{agjō}，最终出于原始印欧语 *\textit{ak-}“\textbf{尖利”}。这个“尖”词根在拉丁一支结出：\textbf{acid}（酸“扎舌头的”）、\textbf{acute}（尖锐的；急性的）、\textbf{acumen}（敏锐）、\textbf{eager}（$<$\textit{acer}“锋利”$\to$热切“像刀刃一样往前冲”）、\textbf{acrid / acrimony}（刻薄）、\textbf{vinegar}（\textit{vinum}+\textit{acre}“酸酒”）、\textbf{oxygen}（\textit{oxy-}“酸”+ \textit{-gen}“生成”，当时误以为氧是酸之母）。一个史前“尖”字，串起了酸、醋、氧与刀锋。
\end{origin}

\vw{emotion}{/ɪˈməʊʃn/}{n.}

\sense{n.}{情感，强烈的感情，激情，情绪}
\col{show / hide / control one's emotions}
\ex{He found it hard to control his emotions at the farewell.}
\col{be filled with / overwhelmed by emotion}
\ex{Overwhelmed by emotion, she could not finish her speech.}
\col{mixed emotions（百感交集）}
\ex{She had mixed emotions about leaving her hometown.}

\begin{gram}{emotion / feeling / mood / passion}
\textbf{feeling} 最宽泛（感觉＋情感）；\textbf{emotion} 较正式，指起伏明显的情绪；\textbf{mood} 心境（好／坏心情 \textit{in a good mood}），与 atmosphere 同源时兼指“氛围”；\textbf{passion} 强烈激情（词源是“受难”，见下框）。形容词链也须辨：\textbf{emotional}（情感的；易激动的）、\textbf{emotive}（煽情的，\textit{an emotive issue} 情绪化议题）。
\end{gram}

\begin{origin}
16 世纪法语 \textit{émotion} $\leftarrow$ 拉丁语 \textit{emovere}（过去分词 \textit{emotus}）“\textbf{移出、搅动而出}”：\textit{e(x)-}（向外）+ \textit{movere}“\textbf{移动”}（原始印欧语 *\textit{meue-}“推动”，日耳曼语表亲为 \textbf{may}？否——但 \textbf{motive, motor, mobile, motion, move} 全族皆出）。原为物理与政治用语“民众的骚动”，“情感”义是 19 世纪浪漫主义时代的产物——把情绪理解为“\textbf{从内心搬动出来}”的东西。同构的 \textbf{promote}（向前搬$\to$促进）、\textbf{remove}（搬走$\to$去除）、\textbf{remote}（搬远了$\to$遥远的）与之共享 \textit{movere}。
\end{origin}

\vw{emphasis}{/ˈemfəsɪs/}{n.\ (pl.\ emphases)}

\sense{n.}{强调，重视；重读}
\col{place / put / lay emphasis on sth（固定搭配）}
\ex{The school places great emphasis on independent learning.}
\col{with emphasis / give special emphasis to}
\ex{She repeated the warning with great emphasis.}
\ex{The report gives special emphasis to early prevention.}
\col{emphatic / emphasize（形容词／动词）}
\ex{The president was emphatic that no deal had been signed.}
\ex{He emphasized the importance of daily practice.}

\begin{gram}{“强调”句式转换：名词与动词的升降级}
名词句 \textit{lay emphasis on} 与动词句 \textit{emphasize} 可互换，但名词句更便于加修饰语升级语气：\textit{lay \textbf{particular / heavy / increasing} emphasis on\dots}。注意 emphasis 复数 emphases（希腊源 -is 规则，详见 crisis 词条）。名词化（nominalization）是学术英语的典型文体特征：动词升级为名词后，句子更“客观”，故说明文与概要写作宜多用。
\end{gram}

\begin{origin}
希腊语 \textit{emphasis}“\textbf{显现、外表、显现于其中之义}”，动词 \textit{emphainein}“使显现”：\textit{en-}（在内）+ \textit{phainein}“\textbf{显现”}（原始印欧语 *\textit{bhā-}“照耀”，同根 \textbf{photo}（\textit{phōs} 光）、\textbf{phenomenon}（显现之物$\to$现象）、\textbf{fancy / fantasy}（\textit{phantasia} 显像）、\textbf{sycophant}（谄媚者“无花果告发者”的显摆者）、\textbf{epiphany}（显现节））。修辞学的 emphasis 原指“\textbf{言外显露之意}”——话说一半而意思自现，后经拉丁语 \textit{emphasis} 进入各现代语言，才成为“加重语气”。
\end{origin}

\vw{employ}{/ɪmˈplɔɪ/}{vt.}

\sense{vt.}{雇用}
\col{employ sb as\dots / be employed in a company}
\ex{The factory employs over two thousand workers.}
\ex{She was employed as a tour guide after graduation.}

\sense{vt.}{运用，使用}
\col{employ a method / technique / tactics（正式用语）}
\ex{The poet employs simple language to convey deep sorrow.}
\ex{The team employed a range of new coaching methods.}
\col{in sb's employ / employee / employer（受雇于人／雇员／雇主）}
\ex{The detective had been in the prince's employ for years.}

\begin{origin}
经由法语 \textit{employer} $\leftarrow$ 拉丁语 \textit{implicare}“\textbf{折入、卷入}”：\textit{im-}（入内）+ \textit{plicare}“折叠”（见 display 词条）。同一个 \textit{implicare} 在不同路径上裂变为三个词：卷入事务 $\to$ \textbf{employ}“使用、雇用”（把人卷进你的事务里）；卷入纠纷 $\to$ \textbf{implicate}（牵连：\textit{implicated in the scandal}）；言外卷有他意 $\to$ \textbf{imply}（暗示）。法语 \textit{employer} 的另一后代是 \textbf{employment}（就业；\textit{employment rate} 就业率），宏观经济学里与 unemployment 相对。
\end{origin}

\vw{encounter}{/ɪnˈkaʊntər/}{vt.\ \&\ n.}

\sense{vt.}{偶然碰到，意外地遇见；遭遇，碰到}
\ex{On their journey they encountered many unexpected difficulties.}
\ex{We encountered an old friend at the airport by chance.}

\sense{n.}{相遇，遭遇，冲突}
\col{an encounter with sb/sth}
\ex{My first encounter with classical music was at primary school.}
\ex{The patrol reported a brief encounter with enemy forces.}
\col{a close encounter（险些相撞；密切接触）}
\ex{Drivers had a close encounter with a falling tree during the storm.}

\begin{origin}
经由古法语 \textit{encontrer}（名词 \textit{encontre}）$\leftarrow$ 晚期拉丁语 \textit{incontra}“\textbf{迎面、相对}”，由 \textit{in-} + \textit{contra}“\textbf{相对、相反”}（见 contradictory 词条）构成——相遇即“\textbf{面对面撞上}”。同根的 \textbf{counter-}（对抗：\textit{counterattack} 反击、\textit{counterclockwise} 逆时针）、\textbf{country}（“对面之地”$\to$身外之地$\to$乡野$\to$国家，语义漂移的奇观）、\textbf{contrary}（相反的）。而 \textbf{encounter} 在医学语篇中已成术语“就诊”（医患的“相遇”），语义仍在生长。
\end{origin}

\vw{enhance}{/ɪnˈhɑːns/}{vt.}

\sense{vt.}{提高，增强}
\col{enhance the quality / value / reputation / appeal}
\ex{The new lighting greatly enhances the appeal of the shop window.}
\ex{Regular practice will enhance your reading speed.}
\col{enhanced security / measures（强化的）}
\ex{Enhanced security measures were introduced after the incident.}

\begin{origin}
盎格鲁-法语 \textit{enhauncer}“\textbf{抬高、举起}”，源自通俗拉丁语 *\textit{inaltiare}：\textit{in-} + 晚期拉丁语 \textit{altiare}“升高”（\textit{altus}“\textbf{高}”，同根 \textbf{altitude} 高度、\textbf{alto} 中音部、\textbf{exalt}（ex- 升起$\to$提升；赞扬）、\textbf{hauteur} 傲慢）。拼写中的 \textbf{h} 是 15 世纪抄写员误仿 \textit{-ance} 类法语词加上的——\textit{enhauncer} 的发音里从来没有 h 的音。故 enhance 字面即“\textbf{拔高}”：质量、价值、声誉的“提高”都是把某物\textbf{举升}一截。
\end{origin}

\vw{enormous}{/ɪˈnɔːməs/}{adj.}

\sense{adj.}{巨大的，庞大的}
\col{an enormous amount / impact / size}
\ex{The Internet has had an enormous impact on daily life.}
\ex{Enormous amounts of money are spent on advertising each year.}
\col{enormously（极大地）}
\ex{The city has changed enormously since the 1990s.}

\begin{gram}{“巨大”的量级与语体}
\textbf{big / large}（中性）；\textbf{enormous}（超出常规，多修饰数量、程度、影响）；\textbf{huge}（口语化，泛）；\textbf{immense}（不可测量的大，语体高：\textit{immense gratitude}）；\textbf{vast}（专长于面积、范围：\textit{vast land}）；\textbf{tremendous}（惊人且常含积极赞叹：\textit{tremendous progress}）；\textbf{giant / colossal}（夸张比喻）。写作建议：数据语境用 enormous，空间语境用 vast，情感语境用 immense / tremendous，避免通篇 huge。
\end{gram}

\begin{origin}
经由拉丁语 \textit{enormis}“\textbf{越出常规的}”：\textit{e(x)-}（超出）+ \textit{norma}“\textbf{木工的直角规、准则”}（词源不详，或与 \textit{gnōrus}“已知”相关；后代有 \textbf{normal} 正常的、\textbf{norm} 标准、\textbf{enormity}（古义“暴行”$\to$今多作“巨大”，曾引发用法之争））。罗马木匠以直角规（norma）划出“正”，越出规矩者即 \textit{enormis}——最初竟带贬义“\textbf{无法无天的}”，18 世纪后才中性化为“大得反常”。同构词 \textbf{abnormal}（ab- 离开规范$\to$反常）与之互补。
\end{origin}

\vw{enthusiastic}{/ɪnˌθjuːziˈæstɪk/}{adj.}

\sense{adj.}{热情的，热心的}
\col{be enthusiastic about sb/sth（固定介词 about）}
\ex{The students were enthusiastic about the science project.}
\ex{He is an enthusiastic supporter of the local football club.}
\col{enthusiasm / enthusiast（名词族）}
\ex{Her enthusiasm for the subject was infectious.}
\ex{He is a model railway enthusiast.}

\begin{origin}
希腊语 \textit{enthousiastēs}“\textbf{神灵附体者}”，出自 \textit{enthousiazein}“受神感应”：\textit{en-}（在内）+ \textit{theos}“\textbf{神}”（原始印欧语 *\textit{dhes-}“神圣力量”，同根 \textbf{theology} 神学、\textbf{atheist} 无神论者、\textbf{pantheon} 泛神庙／万神殿、\textbf{epiphany} 之一义）。柏拉图认为诗人的灵感来自\textbf{神进入其身}（\textit{entheos}“内有神者在”）——热情的本义是“\textbf{体内住进了神}”。17 世纪英语中 enthusiastic 一度特指“宗教狂热者”（带贬义），19 世纪后才泛化为日常的“热忱”。从神启到热心，是一次彻底的世俗化。
\end{origin}

\vw{entitle}{/ɪnˈtaɪtl/}{vt.\ (usually passive)}

\sense{vt.}{给……命名}
\col{be entitled \dots（注意无引号、斜体书名）}
\ex{She wrote a best-selling book entitled \textit{ Wild Swans}.}

\sense{vt.}{使享有权利，使符合资格}
\col{entitle sb to sth / entitle sb to do sth}
\ex{The ticket entitles you to a free drink at the bar.}
\ex{Members are entitled to use the library at any time.}
\col{be entitled to sth（有权享有）}
\ex{Everyone is entitled to their own opinion.}

\begin{gram}{title 的构词族：entitle / subtitle / retitle / titled}
\textbf{en-}（使成）+ \textbf{title} 即“使有标题”；\textbf{subtitle}（sub- 之下$\to$副标题；字幕）、\textbf{entitlement}（应得权益，社会福利语篇高频：\textit{entitlement program}）。口语中 \textbf{entitled} 的新贬义值得关注：\textit{an entitled kid}“\textbf{自以为理应享受}的孩子”——由“被授予权利”滑向“觉得理所当然”，词义正在我们眼前演变。
\end{gram}

\begin{origin}
经由古法语 \textit{entitler} $\leftarrow$ 晚期拉丁语 \textit{intitulare}“\textbf{题名、加封号}”：\textit{in-} + \textit{titulus}“\textbf{铭牌、题签、尊号”}。古代书卷顶端系有小签（\textit{titulus}）标明内容，遂成“标题”；罗马人的凯旋式上高举的标语牌也叫 \textit{titulus}，宣示战功。\textbf{title} 的两个后代路径迥异：一条通往图书馆（书名），一条通往律师事务所（\textbf{title deed} 地契“所有权的名分”、\textbf{title} 冠军头衔）——“名分即权利”的法律思维在 entitle 一词中合二为一。
\end{origin}

\vw{era}{/ˈɪərə/}{n.}

\sense{n.}{时代，年代}
\col{a new / bygone / the postwar era}
\ex{The fall of the Berlin Wall closed one era and opened another.}
\col{the era of\dots}
\ex{We are entering the era of artificial intelligence.}

\begin{origin}
拉丁语 \textit{aera}“\textbf{计数用的铜牌；纪年起点”}，是 \textit{aes}（\textbf{铜、青铜、铜钱}，原始印欧语 *\textit{h₂es-}“铜”）的复数形式。罗马人用小铜牌（\textit{aera}）作筹码计数，历法上遂以某个“\textbf{起点数}”为 \textit{aera}——如西班牙纪元（era hispanica）自公元前 38 年起算。英语 \textbf{era} 于 17 世纪定型为“纪元、时代”；同根词 \textbf{aerate}（充气？不——\textit{aer} 是“空气”，恰是\textbf{另一词源}，勿混：aer$\leftarrow$希腊 \textit{aēr}，era$\leftarrow$拉丁 \textit{aes}，形近而实异，这正是词源学习中“假朋友”的典型）。
\end{origin}

\vw{essence}{/ˈesns/}{n.}

\sense{n.}{本质，精髓}
\col{the essence of sth}
\ex{Competition is the essence of market economy.}
\col{in essence（本质上，书面插入语）}
\ex{In essence, the two plans differ only in scale.}
\col{of the essence（至关重要的）}
\ex{In first aid, speed is of the essence.}

\sense{n.}{香精，精油}
\ex{The shampoo contains natural herbal essences.}

\begin{origin}
中世纪拉丁语 \textit{essentia}“\textbf{存在（之性）}”，由 \textit{esse}“\textbf{是、存在}”派生——\textit{esse} 是拉丁语的不定式“to be”，源自原始印欧语 *\textit{h₁es-}“是”（与英语 \textbf{is}、德语 \textbf{ist}、梵语 \textit{asti} 全部同源，印欧语最有身份的动词）。哲学动词 \textit{esse} 的派生词在英语中俯拾皆是：\textbf{essential}（必不可少的$\to$本质的）、\textbf{absent}（\textit{ab-} 不在$\to$缺席的）、\textbf{present}（\textit{pre-} 在前$\to$在场的；现在的；礼物）、\textbf{represent}（再次呈现）、\textbf{interest}（\textit{inter-} 介于其中“与他事有关”$\to$利益；兴趣）。“本质”就是“使一物\textbf{是其所是}的东西”。
\end{origin}

\vw{essential}{/ɪˈsenʃl/}{adj.}

\sense{adj.}{必不可少的}
\col{be essential to / for sth}
\ex{Vitamin D is essential for healthy bones.}
\col{it is essential that\dots（从句用 should 型虚拟）／it is essential to do sth}
\ex{It is essential that all safety procedures (should) be strictly followed.}
\ex{It is essential to book tickets in advance.}

\sense{adj.}{本质的，基本的}
\col{an essential difference / feature}
\ex{There is an essential difference between the two theories.}
\col{essentials（名词：必需品；要点）}
\ex{Pack only the essentials --- the airline allows one bag.}

\begin{gram}{“it is essential that + (should) do”：形容词引发的虚拟}
虚拟语气不止于 if 从句。\textbf{essential, vital, necessary, important, imperative, crucial, advisable} 等表“必要性／评价”的形容词用于 \textit{It is + adj.\ + that} 结构时，从句谓语用 \textit{(should+) 动词原形}，与 arrange/suggest/demand 等动词后的隐式虚拟（见 arrange 词条）同源，统称“命令性虚拟”。被动省略 should 后形如原形动词，常令读者误认作错误：\textit{It is essential that he \textbf{be} informed.}（$\neq$ is！）。此为改错与语法填空的顶级考点。
\end{gram}

\vw{estimate}{/ˈestɪmeɪt/ vt. \ \ /ˈestɪmət/ n.}{vt.\ \&\ n.}

\sense{vt. /ˈestɪmeɪt/}{估计，估价}
\col{estimate that\dots / estimate sth at\dots}
\ex{Scientists estimate that the reef is over 8,000 years old.}
\ex{The damage was estimated at two million dollars.}
\col{be estimated to do / to be}
\ex{The temple is estimated to have been built in the Tang dynasty.}

\sense{n. /ˈestɪmət/}{估计；估计的成本}
\col{a rough / conservative estimate}
\ex{A rough estimate put the crowd at 50,000.}
\col{by sb's estimate / at a conservative estimate}
\ex{By the committee's estimate, the project will take five years.}

\begin{gram}{名动异读，再遇 estimate}
\textbf{estímate}（动词 /-meɪt/，重音在后）vs.\ \textbf{éstimate}（名词 /-mət/，重音在前），规则同 advocate 词条的“重音转移构词法”。实用技巧：名词 estimate 常与冠词连用（\textit{a rough estimate}）、动词 estimate 常带 that 从句或 at 短语。相关词 \textbf{overestimate / underestimate}（高估／低估）中的 over/under 是前缀而非副词，注意重音一律在 -es-。
\end{gram}

\begin{origin}
经由拉丁语 \textit{aestimare}“\textbf{估定价值}”（词源古老而不确定，一说 \textit{aestimare} $=$ *\textit{ais-temos}“\textbf{切下一块铜}”——罗马人铜块按重量交易，切铜即开价；\textit{aes}“铜”见 era 词条）。其后代遍布罗曼语族：法语 \textit{estimer}（估$\to$敬重：\textit{esteem} 尊敬，对人的“估价”）、意大利 \textit{stimare}。英语 \textbf{estimate}“估价”、\textbf{esteem}“敬重”、\textbf{aim}（$<$\textit{aestimare} 经古法语 \textit{aesmer}“估量着做”$\to$瞄准、打算）竟是三兄弟——瞄准、尊敬与估价，都始于\textbf{掂量一块青铜}。
\end{origin}

\vw{expand}{/ɪkˈspænd/}{vi. \& vt.}

\sense{vi. \& vt.}{增加，扩大；扩展，发展（业务）}
\col{expand rapidly / expand into new markets}
\ex{The chain plans to expand into Southeast Asian markets.}
\ex{The city's population expanded rapidly in the 1990s.}
\col{expanded / expanding（形容词）}
\ex{The airline runs an expanded summer schedule.}

\sense{vi. \& vt.}{详谈，详述}
\col{expand on/upon sth（展开论述）}
\ex{Could you expand on that point? I'm not sure I follow.}

\begin{origin}
拉丁语 \textit{expandere}“\textbf{摊开}”：\textit{ex-}（向外）+ \textit{pandere}“\textbf{铺展”}。\textit{pandere} 的另一复合词 \textit{pandere}（\textit{pan-}?）实出一源二用：一是“张开（门）”，给英语留下 \textbf{expand}、\textbf{expanse}（广阔空间）、\textbf{expansive}（辽阔的；健谈的）；二是“铺开（饼）”，藏在出人意料的地方——\textbf{paste}（面糊“摊平之物”）与 \textbf{pastry}（糕点“摊开的东西”）亦多被归于此根。故 expand 与 pastry 或为远亲：一个是摊开的空间，一个是摊开的面团。同族反义词 \textbf{contract}（con- 一起拉$\to$收缩；合同“拉拢的约定”）与之成对。
\end{origin}

\vw{expectation}{/ˌekspekˈteɪʃn/}{n.}

\sense{n.}{希望，盼望；预料，预期；期望，指望}
\col{expectations of / for sb/sth}
\ex{Parents' expectations of their children are sometimes too high.}
\col{live up to / meet sb's expectations（达到期望）}
\ex{The film failed to live up to audience expectations.}
\col{against / contrary to (all) expectations（出乎意料）}
\ex{Contrary to all expectations, the rookie won the championship.}
\col{beyond (all) expectation（超乎预期）}
\ex{The festival succeeded beyond all expectation.}

\begin{gram}{expect 的三朵姊妹花：expectation / expectancy / life expectancy}
\textbf{expectation} 可数或不可数，“预期、期望”；\textbf{expectancy} 不可数，“期待的状态”（\textit{an air of expectancy} 一派翘首以待的气氛）；\textbf{life expectancy}“预期寿命”是人口学与新闻的高频术语：\textit{Life expectancy has risen to 78.} 相关短语 \textbf{expectant}（期待的；\textit{expectant mothers} 准妈妈们）亦常见于健康类语篇。
\end{gram}

\vw{expose}{/ɪkˈspəʊz/}{vt.}

\sense{vt.}{揭露，揭穿}
\col{expose a scandal / lie / corruption}
\ex{The journalist exposed the factory's illegal waste dumping.}

\sense{vt.}{暴露，显露；使遭受}
\col{be exposed to sth（遭受；接触——高频被动）}
\ex{Children exposed to second-hand smoke are more likely to fall ill.}
\ex{The rocks are exposed to wind and rain all year round.}

\sense{vt.}{使接触}
\col{expose sb to new ideas / cultures}
\ex{Studying abroad exposes students to different ways of thinking.}
\col{exposure（名词：暴露；揭发；曝光）}
\ex{Prolonged exposure to the sun can cause skin cancer.}

\begin{gram}{pose“放置”词族总复习}
拉丁语 \textit{ponere}（过去分词 \textit{positus}）“放置”的家族至此全员到齐：\textbf{com-pose}（放在一起$\to$组成；作曲；使镇静）、\textbf{ex-pose}（放出$\to$暴露；揭露）、\textbf{op-pose}（放在对面$\to$反对）、\textbf{sup-pose}（放在下面作底$\to$假定）、\textbf{im-pose}（放上$\to$强加）、\textbf{de-pose}（拿开$\to$废黜）、\textbf{pro-pose}（向前放$\to$提议；求婚）、\textbf{inter-pose}（放在中间$\to$插话）、\textbf{re-pose}（重放$\to$休息）、\textbf{purpose}（\textit{pro-positum} 向前放置的目标$\to$目的）。凡见 \textit{-pose/-posit-}，先想“放”，词义自明。
\end{gram}

\vw{extraordinary}{/ɪkˈstrɔːdnri/}{adj.}

\sense{adj.}{不寻常的，非凡的}
\col{an extraordinary talent / achievement / sight}
\ex{She showed extraordinary courage for a girl of fourteen.}
\col{extraordinary measures（非常手段）}
\ex{The government took extraordinary measures to contain the outbreak.}

\sense{adj.}{意想不到的，令人惊奇的}
\ex{It's extraordinary that nobody noticed the mistake earlier.}
\col{an extraordinary meeting（特别会议，正式用法）}
\ex{The board called an extraordinary meeting to discuss the merger.}

\begin{gram}{extra- 前缀：在“之外”}
\textbf{extraordinary} $=$ extra（超出）+ ordinary（寻常），“\textbf{超出寻常}”。前缀 extra- 来自拉丁语 \textit{extra}“在外面、超出”，构词力极强：\textbf{extracurricular}（课程之外的$\to$课外的）、\textbf{extraordinary}、\textbf{extravagant}（\textit{vagari} 游荡“游出到界外”$\to$奢侈的）、\textbf{extreme}（最外的$\to$极端的）、\textbf{extraterrestrial}（地球之外的$\to$外星的）。注意它与单词 \textbf{extra}“额外的”同源而分工不同：前者作前缀常译“超／格外”，后者自由使用（\textit{extra pay} 加班费）。
\end{gram}

%================================
\section{F}
%================================

\vw{fascinate}{/ˈfæsɪneɪt/}{vt. \& vi.}

\sense{vt. \& vi.}{深深吸引，迷住}
\col{be fascinated by / with sth（高频被动）}
\ex{As a boy he was fascinated by the stars and kept a telescope by his window.}
\col{fascinating（迷人的）}
\ex{The museum offers a fascinating glimpse into ancient life.}
\col{sth fascinates sb（事物作主语）}
\ex{The way bees communicate has fascinated scientists for decades.}

\begin{origin}
拉丁语 \textit{fascinare}“\textbf{施魔咒、蛊惑}”，源自 \textit{fascinum}——一枚\textbf{罗马人佩戴的护身符}，形为男根状小坠，用来\textbf{反射邪眼}（\textit{oculus malus}）的注视；庞贝古城的门口至今镶着它。动词本指“被邪眼盯住$\to$中邪般不能动弹”，后弱化为“魅力摄人”。1580s 进入英语时仍带巫术色彩（\textit{fascination} 在弥尔顿《失乐园》里是妖术）。“迷人”的词源，竟是一场\textbf{驱邪}。同根词 \textbf{fascist / fascism}（意大利语 \textit{fascio}“束棒”——权力魅惑的另一支，词源上同出“捆束”，已与 fascinate 分道扬镳，故不必强联，但知其别亦佳）。
\end{origin}

\vw{flavour}{/ˈfleɪvər/}{n.\ \&\ vt.}

\sense{n.}{味，味道}
\col{a strong / delicate / smoky flavor}
\ex{The soup has a rich flavor of mushrooms and herbs.}
\col{flavor enhancer / additives（增味剂／添加剂）}
\ex{The label lists no artificial flavors or colors.}

\sense{n.}{特点，特色}
\col{a local / ethnic flavor（地方／民族风味）}
\ex{The old street still keeps some of its original flavor.}

\sense{vt.}{给……调味，加味于}
\col{flavor sth with\dots（注意用 with）}
\ex{She flavored the cake with lemon and vanilla.}

\begin{gram}{五味之外：taste / flavor / seasoning}
\textbf{taste} 泛指味觉与“尝起来”的体验（\textit{It tastes sweet.}），另有“品味”义（\textit{in bad taste} 不得体）；\textbf{flavor} 是风味的整体印象，兼作动词与“特色”义（英式拼 \textbf{flavour}，美式 \textbf{flavor}，本表按英式，但考试从美式时以题目为准）；\textbf{seasoning} 调味料（动词 \textit{season}“加调料”——与“季节”同词而异源？实为同源：\textit{season} $<$ \textit{satio}“使成熟适宜”，蔬菜\textbf{应季}而熟，肉经\textbf{适时}调料而美）。
\end{gram}

\begin{origin}
经由古法语 \textit{flaor／flavor}“\textbf{香气、嗅到的气味}”，多数学者追溯至通俗拉丁语 *\textit{flator}（$\leftarrow$ 拉丁语 \textit{flare}“\textbf{吹、呼气}”，原始印欧语 *\textit{bhle-}“膨胀、吹”：同根 \textbf{inflate}（吹入$\to$通货膨胀的 -flation）、\textbf{conflate}（吹到一起$\to$合并）、\textbf{flavor} 的嗅觉近亲）。最初 \textit{flavor} 主要指\textbf{香气}而非滋味，直至近代才综合为“风味”；\textbf{flavor} $+$ \textbf{savor}（滋味，$<$拉丁 \textit{sapor}，\textbf{sapient}“智者”本义“有品味的人”）共同构成英语的“风味二重奏”。可见英语的“味道”一开始是\textbf{闻}出来的。
\end{origin}

\vw{flow}{/fləʊ/}{vi.\ \&\ n.}

\sense{vi.}{流，流动；充满；涨起}
\col{flow into / out of / through}
\ex{Hundreds of rivers flow into the Yangtze.}
\ex{Wealth flowed into the port city during the boom.}
\col{tears / blood / traffic flows（喻体流动）}
\ex{Traffic flows more smoothly since the new bridge opened.}

\sense{vi.}{涨起（潮水）}
\ex{The tide began to flow, covering the mudflats.}

\sense{n.}{流，流动；连贯；涨潮}
\col{the flow of traffic / goods / ideas}
\ex{China's e-commerce platforms speed up the flow of goods.}
\col{in full flow（滔滔不绝；全速进行）}
\ex{Once the speaker was in full flow, nobody dared interrupt.}
\col{the flow of the essay（行文的连贯）}
\ex{Transition words improve the flow of your writing.}

\begin{origin}
古英语 \textit{flōwan}“流动”，源自原始日耳曼语 *\textit{flōaną}，最终出于原始印欧语 *\textit{plew-}“\textbf{流动、漂浮、航行”}——这个“水”词根的子孙遍布印欧：日耳曼支 \textbf{flow}、\textbf{flood}（洪水）、\textbf{float}（漂）、\textbf{fleet}（舰队“漂着的东西”；疾驰）、\textbf{fly?}（否——fly 是另一词根 *\textit{plew-}“飞”？实为 *\textit{plew-} 的近形异根，恰是词源学的经典陷阱）；拉丁支 \textbf{pluvial}（多雨的，\textit{pluvia} 雨）、\textbf{Pliocene / Pleistocene}（地质纪元“更多雨／更多冰的时代”）；希腊支 \textbf{pneuma}？否——但 \textbf{plasma}（成形之物“流动体”）、\textbf{-pleura}（\textit{pleura} 肋“侧膜”）。一词根，半部水文地理。
\end{origin}

\vw{folk}{/fəʊk/}{adj.\ \&\ n.}

\sense{adj.}{民间的，民俗的}
\col{folk music / dance / art / tales}
\ex{The festival features folk dances from a dozen ethnic groups.}
\ex{Folk tales were passed down orally for generations.}

\sense{n.}{民间音乐}
\ex{He developed a love of folk while at college.}

\sense{n.}{人们；各位；家人}
\col{folks（口语：各位，朋友们）}
\ex{That's all for now, folks --- thanks for listening.}
\col{sb's folks（某人的家人，口语）}
\ex{I'm going home to see my folks this weekend.}
\col{old folks（老人们）；country folk（乡里人）}
\ex{The old folks still gather under the banyan tree on summer evenings.}

\begin{gram}{folk 在学术英语中的妙用：“民俗”术语族}
语言学与文化研究高频术语皆由 folk 构成：\textbf{folklore}（民俗学；民间传说）、\textbf{folktale}（民间故事，区别于 legend 传说与 myth 神话）、\textbf{folk art / folk custom}（民间艺术／习俗）、\textbf{folk medicine}（民间医药）、\textbf{folk etymology}（\textbf{民俗词源}——民众误按熟词重新解析词形，如 \textit{sparrowgrass}$<$asparagus、\textit{bridegroom} 中的 -groom 本为古英语 \textit{guma}“人”，被误解为“马夫”；本讲义词源栏正是它的解毒剂）、\textbf{wisdom folk}？作 \textbf{folk wisdom}（民间智慧）。
\end{gram}

\begin{origin}
古英语 \textit{folc}“\textbf{民众、部队”}（原始日耳曼语 *\textit{fulka}，同源德语 \textbf{Volk}“人民、民族”——19 世纪浪漫主义与政治史中分量极重的词；荷兰语 \textit{volk}；古诺斯语 \textit{fólk} 兼指“战士队”）。印欧词源一说与 *\textit{pelh\u{1}-}“填满”相关（“众多填满之地”），故与 \textbf{full}、\textbf{populate}（\textit{populus} 人民）遥为表亲。英语中 folk 保留了古拙的乡野气息，正因如此被 19 世纪学者选中组成 \textbf{folklore}（1846 年 W.\,J.\,Thoms 仿德语 \textit{Volkskunde} 所造）——一个古老的词，命名了一门年轻的学科。
\end{origin}

\vp{for good measure}

\sense{adv. 短语}{作为额外增添}
\ex{The shop threw in a free pen for good measure.}
\ex{He told the joke again, adding a new punchline for good measure.}

\begin{gram}{measure 的“量度”短语族}
\textbf{for good measure}（额外再加）原是售货用语“\textbf{分量之外再添一点}”；同族的：\textbf{a measure of}（一定程度的：\textit{a measure of success}）、\textbf{beyond measure}（无可估量：\textit{grateful beyond measure}）、\textbf{in a / some measure}（在某种程度上，书面）、\textbf{made to measure}（量体裁衣的）、\textbf{take measures / steps to do}（采取措施——measure 此处已是“手段”）、\textbf{half measure}（半吊子措施，多用于否定 \textit{no half measures}）。一条“量”的词义链：尺寸$\to$程度$\to$手段。
\end{gram}

\end{document}
